 In the following, we discuss similarities, drawbacks and benefits of Neural SDDDE compared with the following models: NODE, ANODE and Neural Laplace. Compared to these techniques, we recall that Neural SDDDE is a direct method for solving DDEs, specifically designed to handle delayed systems; it does not require lifting of the problem to higher-dimensional spaces, with the benefit of preserving physical interpretability. This is not the case for ANODE, the augmented version of ODE, where flexibility is obtained by introducing an higher dimentional space. Lastly, Neural Laplace can solve a broader class of differential equations, although with some limitations that we pinpoint in the following.

\subsection{Model comparison} From the theoretical viewpoint, the Laplace transformation is often a tool used in proofs on DDEs \citep{Cook_book_ref_dde} since it allows to transform linear functional equations in $\mathbf{f}(\mathbf{x}(t))$ involving derivative and differences into linear equations involving only $\mathbf{F}(\mathbf{s})$. Thus, time-dependent and constant delay DDEs are transformed into linear equations of $\mathbf{F}(\mathbf{s})$ using the Laplace transform. This transformation enables Neural Laplace to bypass the explicit definition of delays, whereas Neural SDDDE needs the delays to be specified unless the vector flow $\mathbf{f}$ and delays are learnt jointly. These observations tie Neural Laplace to Neural SDDDE as they can be seen as similar models but living in different domains. However, a limitation of Laplace transformation-based approaches is that they are not defined for DDEs with state-dependent delays, thus limiting the class of time-delayed equation that one can solve with this technique. 

Neural Laplace is a model that needs memory initialized latent variables, i.e., a long portion of the trajectory needs to be fed in order to get a reasonable representation of the latent variable. In contrast, by design, NODE and ANODE require information only at the initial time to predict the system dynamics. From this viewpoint, Neural SDDDE is somewhat a hybrid approach since a history function $\phi(t)$ is required. $\phi$ needs to be provided for $t\in[-\tau_{\mathrm{max}},0]$, where $\tau_{\mathrm{max}}$ is the maximum delay encountered during integration. This is more restricting than solely relying on an initial condition $\mathbf{x}(0)$ but far less than memory-based latent variables methods. Moreover, the training and testing schemes for Neural Laplace is constrained by this very same observation since future events can only be predicted after a certain observation time. This also makes the length of the trajectory to feed to Neural Laplace a hyperparameter to tune. This is not the case with NODE, ANODE and our approach. 

\subsection{Memory and computational complexity} Neural SDDDEs rely on an ODE solver, thus function evaluations are associated with the same computational cost as NODEs. However, Neural SDDDE has some extra constraints making the method more computationally involved. Hereafter, we compare the complexity of these two schemes; we define $S$ as the number of stages in the Runge-Kutta (RK) scheme used for the time-integration, $N$ the total number of integration steps, and $d$ the state's dimension. 

\paragraph{Memory complexity} DDE integration necessitates keeping a record of all previous states in memory due to the presence of delayed terms. This is because at any given time $t$, the DDE solver must be able to accurately determine $\mathbf{x}(t-\tau)$ through interpolation of the stored state history. This extra amount of extra memory needed depends on the solver used. For example, the additional memory required when using a RK solver for one trajectory is $O(SNd)$. The model's memory footprint is not affected by the number of delays. 

\paragraph{Time complexity} In comparison to NODE, the solution estimate $\widehat{\mathbf{x}}$ needs to be evaluated for each delayed state argument, i.e., $\mathbf{x}(t-\tau_i)$. The cost of evaluating the interpolant is small compared to the cost of computing its coefficients. Similarly, the time complexity is conditioned by the solver used. For example, for a RK scheme, the coefficient computation scales linearly with the number of stages. Hence, Neural SDDDE adds a time cost $O(SD\,d)$ compared to NODE for each vector field function evaluation.